\documentclass[11pt]{article}

\usepackage[T1]{fontenc}
\usepackage{amsmath,amssymb,amsthm}
\usepackage[margin=1in]{geometry}
\usepackage[colorlinks=true,linkcolor=blue,citecolor=blue,urlcolor=blue]{hyperref}

\newtheorem{theorem}{Theorem}
\newtheorem{proposition}[theorem]{Proposition}
\newtheorem{lemma}[theorem]{Lemma}
\newtheorem{corollary}[theorem]{Corollary}
\newtheorem{conjecture}[theorem]{Conjecture}
\theoremstyle{remark}
\newtheorem{remark}[theorem]{Remark}

\newcommand{\R}{\mathbb R}
\newcommand{\norm}[1]{\left\|#1\right\|}
\newcommand{\set}[1]{\left\{#1\right\}}
\newcommand{\cZ}{\mathfrak Z}
\newcommand{\cL}{\mathcal L}
\newcommand{\esssup}{\operatorname*{ess\,sup}}
\setlength{\emergencystretch}{2em}

\title{The q4 Adjoint Dissipation Is Pinned to
High-Amplitude Drift Slabs}
\author{John Lawson and Codex}
\date{25 July 2026}

\begin{document}
\maketitle

\begin{abstract}
This round audits reverse-time Caccioppoli and Meyers improvement for the
conditional exact q4 projected-Oseen entrance.  The exact logarithmic
weak-\(L^3\) clock yields a Lorentz--Zygmund closing norm for the
adjoint gradient that is strictly weaker than the previously recorded
\(L_t^{22/9,2}L_x^2\) condition.  If the entrance survives, this norm
must diverge on every terminal interval.

Equivalently, finite unweighted adjoint dissipation is forced onto
arbitrarily high drift-amplitude slabs with a divergent first moment.
The resulting pointwise gradient spikes have the exact
\(M^{9/4}\) power up to arbitrarily slow summable losses.

Differentiated energy, current BMO-skew reverse H\"older theory, and
generic parabolic smoothing do not close this endpoint.  The direct
enstrophy hull always carries an unavailable fourth power of primal
enstrophy; existing reverse H\"older exponents depend on a uniform BMO
norm which diverges across q4 record blocks; projected pressure is known
only at the base forcing exponent; and a heat-flow construction rules out
every generic \(L_t^{2+\epsilon}L_x^2\) gradient gain from \(L^2\)
boundary data.  No Clay alternative is proved.
\end{abstract}

\section{Epistemic status and exact input}

Every repository theorem below is conditional on the energy-efficient
exact q4 branch and its projected-Oseen entrance.  Those antecedents are
not derived for arbitrary Clay data.  External mathematical review is
pending.

On every terminal interval \(I=(t_0,T)\), put
\begin{equation}
\label{eq:MY}
 M(t):=\norm{u(t)}_{L^{3,\infty}},
 \qquad
 Y(t):=\norm{\nabla a(t)}_2,
 \qquad
 p_*:=\frac{11}{2}.
\end{equation}
The exact first-record clock gives, for large \(\lambda\),
\begin{equation}
\label{eq:clock}
 \left|\set{t\in I:M(t)>\lambda}\right|
 \lesssim
 \frac{\lambda^{-p_*}}{\log(e+\lambda)}.
\end{equation}
The preceding weighted-renormalisation theorem gives
\begin{equation}
\label{eq:weighted-infinite}
 \boxed{
 \int_I M(t)Y(t)^2\,dt=\infty,
 }
\end{equation}
although
\begin{equation}
\label{eq:unweighted-finite}
 \int_IY(t)^2\,dt<\infty.
\end{equation}

The Clay problem remains unsolved.

\section{The logarithmically relaxed closing norm}

Let \(M_I^*\) and \(Y_I^*\) denote decreasing rearrangements on \(I\).
For \(0<s<|I|\), set
\begin{equation}
\label{eq:ell}
 \ell_I(s):=\log\left(e+\frac{|I|}{s}\right)
\end{equation}
and define
\begin{equation}
\label{eq:Znorm}
 \boxed{
 \cZ_I(Y)^2
 :=
 \int_0^{|I|}
 s^{-2/11}\ell_I(s)^{-2/11}
 \bigl(Y_I^*(s)\bigr)^2\,ds.
 }
\end{equation}

\begin{lemma}[Clock inversion]
\label{lem:clock-inversion}
There is a constant \(C_I\) such that
\begin{equation}
\label{eq:Mstar}
 M_I^*(s)
 \leq
 C_I\left[
 1+s^{-2/11}\ell_I(s)^{-2/11}
 \right].
\end{equation}
\end{lemma}

\begin{proof}
Only small \(s\) matters.  For a constant \(A>1\), define
\[
 \lambda_s
 :=
 A s^{-1/p_*}\ell_I(s)^{-1/p_*}.
\]
For sufficiently small \(s\),
\[
 \log(e+\lambda_s)\asymp\ell_I(s).
\]
Consequently \eqref{eq:clock} gives
\[
 \left|\set{M>\lambda_s}\right|
 \lesssim A^{-p_*}s.
\]
Choose \(A\) so that the last expression is at most \(s\).  The
definition of decreasing rearrangement gives
\(M_I^*(s)\leq\lambda_s\).  A constant controls the remaining compact
range of \(s\).
\end{proof}

\begin{theorem}[Lorentz--Zygmund closure]
\label{thm:Zclosure}
If
\[
 \cZ_I(Y)<\infty,
\]
then
\[
 \int_I MY^2<\infty.
\]
Hence every q4 entrance satisfies
\begin{equation}
\label{eq:Zdivergence}
 \boxed{\cZ_I(Y)=\infty}
\end{equation}
on every terminal interval.
\end{theorem}

\begin{proof}
Hardy--Littlewood rearrangement and
Lemma~\ref{lem:clock-inversion} give
\begin{align*}
 \int_I MY^2
 &\leq
 \int_0^{|I|}
 M_I^*(s)\bigl(Y_I^*(s)\bigr)^2\,ds\\
 &\lesssim_I
 \int_IY^2+\cZ_I(Y)^2.
\end{align*}
Both terms on the last line are finite under the hypotheses.  This
contradicts \eqref{eq:weighted-infinite} for an entrance.
\end{proof}

\begin{corollary}[Relation to the previous target]
\label{cor:old-target}
The condition
\[
 Y\in L^{22/9,2}(I)
\]
implies \(\cZ_I(Y)<\infty\), but the converse is false.
\end{corollary}

\begin{proof}
The Lorentz norm satisfies
\[
 \norm{Y}_{L^{22/9,2}}^2
 \asymp
 \int_0^{|I|}
 s^{9/11-1}\bigl(Y_I^*(s)\bigr)^2\,ds
 =
 \int_0^{|I|}
 s^{-2/11}\bigl(Y_I^*(s)\bigr)^2\,ds.
\]
This dominates \eqref{eq:Znorm}.  For strictness, near \(s=0\) take
\[
 \bigl(Y_I^*(s)\bigr)^2
 =
 s^{-9/11}\bigl[\log(e/s)\bigr]^{-10/11}.
\]
The logarithm in \eqref{eq:Znorm} makes its integral convergent, while
the displayed Lorentz integral diverges.
\end{proof}

\subsection{Strong \texorpdfstring{\(L^{22/9}\)}{L22/9} does not suffice}

\begin{proposition}[Endpoint scalar saturation]
\label{prop:scalar-saturation}
The clock distribution \eqref{eq:clock} and
\(Y\in L^{22/9}(I)\) do not, by themselves, imply
\(\int_I MY^2<\infty\).
\end{proposition}

\begin{proof}
Work on \(0<s<e^{-e}\) and write
\[
 L(s):=\log(e/s),
 \qquad
 L_2(s):=\log L(s).
\]
Define decreasing functions
\[
 M_0(s):=s^{-2/11}L(s)^{-2/11}
\]
and
\[
 Z_0(s)
 :=
 s^{-9/11}L(s)^{-9/11}L_2(s)^{-1},
 \qquad
 Y_0:=Z_0^{1/2}.
\]
Monotone inversion shows that \(M_0\) has distribution
\(\lambda^{-11/2}/\log\lambda\), up to constants.  Moreover,
\[
 \int_0^{e^{-e}}Y_0^{22/9}\,ds
 =
 \int_0^{e^{-e}}
 \frac{ds}{sL(s)L_2(s)^{11/9}}
 <\infty.
\]
On the other hand,
\[
 \int_0^{e^{-e}}M_0Y_0^2\,ds
 =
 \int_0^{e^{-e}}
 \frac{ds}{sL(s)L_2(s)}
 =\infty.
\]
This is a scalar rearrangement ledger, not a Navier--Stokes example.
\end{proof}

\section{The amplitude-slab pincer}

For large integers \(n\), set
\begin{equation}
\label{eq:slabs}
 A_n:=\set{t\in I:2^n<M(t)\leq2^{n+1}},
 \qquad
 d_n:=\int_{A_n}Y(t)^2\,dt.
\end{equation}

\begin{theorem}[Forced high-amplitude dissipation]
\label{thm:slab}
The slab ledger satisfies
\begin{equation}
\label{eq:slab-ledger}
 \sum_nd_n<\infty,
 \qquad
 \boxed{
 \sum_{n\geq N}2^nd_n=\infty
 \quad\text{for every }N.
 }
\end{equation}
For every positive summable sequence \((\varepsilon_n)\), infinitely
many \(n\) satisfy
\begin{equation}
\label{eq:dn-floor}
 d_n\geq2^{-n}\varepsilon_n.
\end{equation}
For each such \(n\), some Lebesgue point \(t_n\in A_n\) satisfies
\begin{equation}
\label{eq:Y-spike}
 \boxed{
 Y(t_n)
 \gtrsim
 2^{9n/4}(n\varepsilon_n)^{1/2}.
 }
\end{equation}
\end{theorem}

\begin{proof}
The first sum in \eqref{eq:slab-ledger} is bounded by
\eqref{eq:unweighted-finite}.  Since \(M\leq2^{n+1}\) on \(A_n\),
\[
 \int_I MY^2
 \leq
 C_N\int_IY^2
 +2\sum_{n\geq N}2^nd_n.
\]
Equation \eqref{eq:weighted-infinite} forces every tail series to
diverge.

If \eqref{eq:dn-floor} failed eventually, that tail would be bounded by
\(\sum_n\varepsilon_n\).  Finally, \eqref{eq:clock} gives
\[
 |A_n|
 \lesssim
 \frac{2^{-11n/2}}{n}.
\]
At a Lebesgue point above a fixed fraction of the slab average,
\[
 Y(t_n)^2
 \gtrsim
 \frac{d_n}{|A_n|}
 \gtrsim
 n\varepsilon_n2^{9n/2},
\]
which proves \eqref{eq:Y-spike}.  The selected indices tend to infinity.
Smoothness on every compact preterminal interval makes \(M\) locally
bounded, so the corresponding times approach \(T\); after passing to a
subsequence they increase to \(T\).
\end{proof}

\begin{corollary}[A concrete near-\texorpdfstring{\(M^{9/4}\)}{M9/4}
spike]
\label{cor:concrete-spike}
There are times \(t_n\uparrow T\) such that
\begin{equation}
\label{eq:concrete-spike}
 \boxed{
 Y(t_n)
 \gtrsim
 \frac{M(t_n)^{9/4}}
 {\log\log(e^e+M(t_n))}.
 }
\end{equation}
\end{corollary}

\begin{proof}
Use
\(\varepsilon_n=[n(\log n)^2]^{-1}\) in
Theorem~\ref{thm:slab}.  This sequence is summable, and
\(M(t_n)\asymp2^n\) on \(A_n\).
\end{proof}

\begin{corollary}[Pointwise amplitude barrier]
\label{cor:pointwise-barrier}
For every \(\beta>0\), \(C>0\), and terminal interval \(I\), the bound
\begin{equation}
\label{eq:pointwise-ceiling}
 Y(t)^2
 \leq
 C\left[
 1+\frac{M(t)^{9/2}}{\log(e+M(t))^\beta}
 \right]
\end{equation}
cannot hold almost everywhere.
\end{corollary}

\begin{proof}
Layer cake and \eqref{eq:clock} give
\[
 \int_I
 \frac{M^{11/2}}{\log(e+M)^\beta}\,dt
 <\infty,
\]
because the high-amplitude integral is bounded by
\[
 C\int^\infty
 \frac{d\lambda}
 {\lambda[\log(e+\lambda)]^{1+\beta}}.
\]
Multiplying \eqref{eq:pointwise-ceiling} by \(M\) would contradict
\eqref{eq:weighted-infinite}.
\end{proof}

\section{Differentiated-energy obstruction}

In reverse time, write
\begin{equation}
\label{eq:reverse}
 \partial_\tau v-\nu\Delta v
 +(b\cdot\nabla)v+\nabla Q=0,
 \qquad
 \nabla\cdot b=\nabla\cdot v=0.
\end{equation}
Set
\[
 X:=\norm{\nabla b}_2,
 \qquad
 Y:=\norm{\nabla v}_2,
 \qquad
 Z:=\norm{\Delta v}_2.
\]

\begin{proposition}[The enstrophy hull retains \(X^4\)]
\label{prop:X4}
At smooth positive reverse times,
\begin{equation}
\label{eq:X4-basic}
 \frac12\frac d{d\tau}Y^2
 +\frac{\nu}{2}Z^2
 \lesssim
 \nu^{-3}X^4Y^2.
\end{equation}
More generally, for \(0<\alpha\leq1\), let
\[
 \frac1{r_\alpha}:=\frac13-\frac{\alpha}{6}.
\]
The interpolation
\[
 \norm{b}_{r_\alpha}
 \lesssim M^{1-\alpha}X^\alpha
\]
produces the Gr\"onwall coefficient
\begin{equation}
\label{eq:full-X4}
 \boxed{
 M^{4/\alpha-4}X^4.
 }
\end{equation}
\end{proposition}

\begin{proof}
Pair \eqref{eq:reverse} with \(-\Delta v\).  Pressure vanishes and
\[
 \int(b\cdot\nabla)v\cdot(-\Delta v)
 =
 \int
 \partial_kb_\ell\,\partial_\ell v_i\,\partial_kv_i.
\]
Therefore
\[
 |\text{transport}|
 \leq
 X\norm{\nabla v}_4^2
 \lesssim
 XY^{1/2}Z^{3/2}.
\]
Young's inequality proves \eqref{eq:X4-basic}.

For the full hull, choose \(q_\alpha\) by
\[
 \frac1{r_\alpha}+\frac1{q_\alpha}+\frac12=1.
\]
Then
\[
 \norm{\nabla v}_{q_\alpha}
 \lesssim
 Y^{1-\gamma_\alpha}Z^{\gamma_\alpha},
 \qquad
 \gamma_\alpha:=\frac3{r_\alpha}=1-\frac{\alpha}{2}.
\]
The transport work is at most
\[
 M^{1-\alpha}X^\alpha
 Y^{1-\gamma_\alpha}Z^{1+\gamma_\alpha}.
\]
The Young exponent on the coefficient is
\[
 \frac2{1-\gamma_\alpha}=\frac4\alpha.
\]
This yields
\[
 \left(M^{1-\alpha}X^\alpha\right)^{4/\alpha}Y^2
 =
 M^{4/\alpha-4}X^4Y^2.
\]
\end{proof}

\begin{remark}
The q4 energy ledger gives \(X\in L^2_\tau\), not \(X\in L^4_\tau\).
At the spatial endpoint \(r=3\), one instead obtains
\[
 |\text{transport}|\lesssim MZ^2,
\]
which is absorbable only for uniformly small \(M/\nu\).  The terminal
record amplitudes diverge.  Thus direct differentiated energy does not
reach Theorem~\ref{thm:Zclosure}.
\end{remark}

\section{BMO-skew reverse H\"older audit}

\begin{lemma}[Weak-\(L^3\) drift has a skew BMO potential]
\label{lem:BMO-potential}
For a divergence-free \(b\in L^{3,\infty}(\R^3)\), define
\begin{equation}
\label{eq:D}
 D_{ij}
 :=
 \partial_j(-\Delta)^{-1}b_i
 -\partial_i(-\Delta)^{-1}b_j.
\end{equation}
Then \(D\) is skew,
\[
 \partial_iD_{ij}=b_j,
\]
and
\begin{equation}
\label{eq:BMO-bound}
 \norm{D}_{\mathrm{BMO}}
 \lesssim
 \norm{b}_{L^{3,\infty}}.
\end{equation}
\end{lemma}

\begin{proof}
The algebraic statements follow from
\(\nabla\cdot b=0\).  The kernel in \eqref{eq:D} has order \(-1\).
For a ball \(B=B(x_0,r)\), split \(b\) into \(2B\) and its complement.
The near contribution has mean oscillation at most
\[
 Cr^{-2}\int_{2B}|b|
 \lesssim
 \norm{b}_{L^{3,\infty}}.
\]
For the far contribution, subtract its value at \(x_0\).  On the
\(k\)-th dyadic annulus the kernel difference is
\(O(r/(2^kr)^3)\), while weak-\(L^3\) gives an \(L^1\) mass at most
\(C\norm{b}_{L^{3,\infty}}(2^kr)^2\).  The resulting
\(C2^{-k}\norm{b}_{L^{3,\infty}}\) series is summable.
\end{proof}

The pressure-free component equation can consequently be placed in
divergence form with symmetric part \(\nu I\) and skew part \(D\).
Zhang's current arXiv v5 preprint~\cite{ZhangBMO} claims local
Caccioppoli and gradient reverse H\"older estimates for parabolic
systems whose real block-diagonal skew part lies in
\(L^\infty_t\mathrm{BMO}_x\).  It yields some \(p>2\) sufficiently
close to two; the exponent and constants depend on the uniform BMO
norm.

\begin{proposition}[Why the preprint estimate does not close q4]
\label{prop:BMO-mismatch}
The cited local reverse H\"older claim does not imply
\(\cZ_I(Y)<\infty\) from the current q4 ledger.
\end{proposition}

\begin{proof}
On a first-record block with \(M\leq K\),
Lemma~\ref{lem:BMO-potential} gives a uniform BMO norm \(O(K)\).
Across terminal blocks \(K\to\infty\), so the source supplies no
record-uniform exponent above two.

The result is local in space-time, whereas \eqref{eq:Znorm} is a
terminal mixed norm of the global \(L_x^2\) gradient.  Finally, the
projected equation is
\begin{equation}
\label{eq:pressure-forcing}
 \partial_\tau v
 -\nabla\cdot((\nu I-D)\nabla v)
 =
 -\nabla Q
 =
 \nabla\cdot(-QI).
\end{equation}
Calder\'on--Zygmund and Lorentz Sobolev estimates give
\begin{equation}
\label{eq:Qbase}
 \norm{Q(\tau)}_2
 \lesssim
 M(\tau)Y(\tau).
\end{equation}
Thus \(Q\in L^2_{\tau,x}\) on a block with \(M\leq K\).  The claimed
\(p>2\) conclusion requires the divergence forcing in the corresponding
\(L^p\) class.  Requiring \(Q\in L^p\) is already a
higher-integrability input and is circular here.
\end{proof}

\section{Heat flow forbids a generic temporal gain}

\begin{proposition}[Simultaneous heat-flow saturation]
\label{prop:heat}
There is a divergence-free \(f\in L^2(\R^3)\) such that
\[
 v(\tau):=e^{\nu\tau\Delta}f
\]
satisfies
\begin{equation}
\label{eq:heat-energy}
 \int_0^\infty\norm{\nabla v(\tau)}_2^2\,d\tau<\infty,
\end{equation}
but for every \(p>2\),
\begin{equation}
\label{eq:heat-no-p}
 \int_0^1\norm{\nabla v(\tau)}_2^p\,d\tau=\infty.
\end{equation}
\end{proposition}

\begin{proof}
Choose mutually disjoint solenoidal Fourier annuli at frequencies
\(N_k=4^k\), and choose \(f_k\) in those annuli with
\(\norm{f_k}_2=k^{-1}\).  Orthogonality gives
\[
 f:=\sum_{k\geq1}f_k\in L^2.
\]
Plancherel and Tonelli give the exact energy integral
\[
 \int_0^\infty
 \norm{\nabla e^{\nu\tau\Delta}f}_2^2\,d\tau
 =
 \frac1{2\nu}\norm{f}_2^2.
\]

There are disjoint time intervals \(J_k\), each of length comparable to
\(N_k^{-2}\), on which the heat multiplier is bounded below on the
\(k\)-th annulus.  Hence
\[
 \norm{\nabla e^{\nu\tau\Delta}f_k}_2
 \gtrsim \frac{N_k}{k}
 \qquad(\tau\in J_k).
\]
For \(p>2\), orthogonality and disjointness yield
\[
 \begin{aligned}
 \int_0^1\norm{\nabla e^{\nu\tau\Delta}f}_2^p\,d\tau
 &\geq
 \sum_k\int_{J_k}
 \norm{\nabla e^{\nu\tau\Delta}f_k}_2^p\,d\tau\\
 &\gtrsim
 \sum_kN_k^{p-2}k^{-p}
 =\infty.
 \end{aligned}
\]
\end{proof}

\begin{remark}
Taking \(b=0\) makes the drift a smooth self-generated
Navier--Stokes trajectory and gives \(Q=0\).  The example has the
nonzero strong trace \(f\), not the q4 zero weak trace.  It proves that
parabolic smoothing and self-generation of the coefficient alone cannot
give the desired temporal exponent.  A successful theorem must use the
zero weak trace together with the nonzero same-trajectory cross defect.
\end{remark}

\section{What changed in this round}

\subsection*{Formal findings, subject to external review}

\begin{enumerate}
\item The Lorentz--Zygmund norm \eqref{eq:Znorm}, strictly weaker than
the previous \(L_t^{22/9,2}\) target, closes q4 conditionally.
\item Every survivor has the divergent amplitude ledger
\eqref{eq:slab-ledger} and the near-\(M^{9/4}\) spikes
\eqref{eq:Y-spike}--\eqref{eq:concrete-spike}.
\item Strong \(L_t^{22/9}\) alone is insufficient under the clock
distribution, by Proposition~\ref{prop:scalar-saturation}.
\item The full direct differentiated-energy hull retains the unavailable
\(X^4\) factor \eqref{eq:full-X4}.
\item The pinned BMO-skew reverse H\"older preprint does not cross the
unbounded-record, mixed-norm, or projected-pressure gates.
\item Even heat flow has no generic temporal
\(L_t^{2+\epsilon}L_x^2\) gradient gain at an \(L^2\) boundary.
\end{enumerate}

\subsection*{Closed shortcuts}

\begin{enumerate}
\item Generic Caccioppoli or heat smoothing cannot prove the target.
\item The BMO-skew reverse H\"older exponent cannot be treated as
uniform across q4 record blocks.
\item The base estimate \(Q\in L^2\) cannot be used as the \(L^p\),
\(p>2\), forcing required to prove the same \(p\)-gain.
\item Strong \(L^{22/9}\) cannot replace the exact rearrangement norm.
\item The forced lower spikes are not upper estimates or contradictions
by themselves.
\end{enumerate}

\subsection*{Open obligations}

\begin{enumerate}
\item Prove \(\cZ_I(Y)<\infty\) from the same-trajectory cross defect.
\item Equivalently, prove a summable amplitude-slab anti-correlation.
\item Develop a pressure-compatible boundary reverse H\"older estimate
whose gain is uniform across the unbounded first-record coefficients,
or obtain the missing pressure equiintegrability from the cross defect.
\item Alternatively close the Fourier-commutator or
\(\mathcal H^1\)-capacity routes from the preceding round.
\item Treat slower clocks, divergent normalised energy, and the other
Clay alternatives separately.
\end{enumerate}

\begin{conjecture}[Same-trajectory Lorentz--Zygmund gain]
Under the exact q4 hypotheses, on some terminal interval \(I\),
\[
 \cZ_I\bigl(\norm{\nabla a}_2\bigr)<\infty.
\]
\end{conjecture}

Theorem~\ref{thm:Zclosure} would then exclude the q4 terminal defect.
The conjecture is not proved.

\section*{Source provenance}

The pinned arXiv v5 source archive for
\cite{ZhangBMO} has SHA-256
\[
\texttt{98ba2c1b92ceff3946f5fb67e2ec992632c5f523da5861f1966fd01e467d17b1}.
\]
Its exact file \texttt{Gz967sa.tex} has SHA-256
\[
\texttt{02f6f3c055992464d99adf0ade44bd4e608bef768a85d0cd19dd5b84de12755c}.
\]
Source lines 367--381 state the uniform
\(L^\infty_t\mathrm{BMO}_x\) assumption, the local \(p>2\)
conclusions, and norm dependence.  Lines 607--630 give the improved
Caccioppoli and reverse H\"older estimates.  Lines 717--719 perform the
Gehring gain.  This is recorded as a preprint claim, not as an
established peer-reviewed theorem.

\begin{thebibliography}{9}
\raggedright

\bibitem{ZhangBMO}
G.~Zhang,
\emph{On regularity for nonhomogeneous parabolic systems with a
skew-symmetric part in BMO},
arXiv:2403.01741v5,
\href{https://arxiv.org/abs/2403.01741}
{arXiv:2403.01741}.

\end{thebibliography}

\end{document}
